\section{Preliminaries} \label{Prel}

\subsection{Dirac system and dNLS}

We denote by $u$ the fundamental solution of system~\eqref{1.5} normalized by
the condition
\begin{gather*}  % \label{nc}
u(0,z)=I_m, \qquad m=m_1+m_2.
\end{gather*}
\begin{Definition} \label{CyW2} Let Dirac system~\eqref{1.5} on $[0, \infty)$
 be given and assume that $V$ is locally summable.
Then the Weyl function $\vp$ is an $m_2 \times m_1$ holomorphic matrix function, which satisf\/ies the inequality
\begin{gather}  \label{2.0}
\int_0^{\infty}
\begin{bmatrix}
I_{m_1} & \vp(z)^*
\end{bmatrix}
u(x,z)^*u(x,z)
\begin{bmatrix}
I_{m_1} \\ \vp(z)
\end{bmatrix}dx< \infty .
\end{gather}
\end{Definition}

The following proposition is proved in \cite{FKRSp1} (see also \cite[Section~2.2]{SaSaR}).
\begin{Proposition} The Weyl function always exists and it is unique.
\end{Proposition}

In order to construct the Weyl function, we introduce a class of $m \times m_1$ matrix functions
$\clp(z)$, which are an immediate analog of the classical pairs
of parameter matrix functions. Namely, the matrix functions
$\clp(z)$ are meromorphic in $\BC_+$ and satisfy
(excluding, possibly, a discrete set of points)
the following relations
\begin{gather}\label{2.1}
\clp(z)^*\clp(z) >0, \qquad \clp(z)^* j \clp(z) \geq 0, \qquad z\in \BC_+.
\end{gather}
It is said that $\clp(z)$ are nonsingular (i.e., the f\/irst inequality in~\eqref{2.1} holds)
and with property-$j$ (i.e., the second inequality in~\eqref{2.1} is valid).
Relations~\eqref{2.1} imply (see, e.g., \cite{FKRSp1}) that
\begin{gather*}  %\label{wd}
\det \big(\begin{bmatrix}
I_{m_1} & 0
\end{bmatrix}u(x,z)^{-1}\clp(z)\big)\not= 0.
\end{gather*}
\begin{Definition} \label{set}
The set $\cln(x,z)$ of M\"obius transformations is the set of values $($at the f\/ixed $x \in [0, \infty)$, $z\in \BC_+)$
of matrix functions
\begin{gather*}%\label{2.2}
\vp(x,z,\clp)=\begin{bmatrix}
0 &I_{m_2}
\end{bmatrix}u(x,z)^{-1}\clp(z)\big(\begin{bmatrix}
I_{m_1} & 0
\end{bmatrix}u(x,z)^{-1}\clp(z)\big)^{-1},
\end{gather*}
where $\clp(z)$ are nonsingular matrix functions
 with property-$j$.
 \end{Definition}

 \begin{Remark}\label{MB}
 It was shown in \cite{FKRSp1} that a family $\cln(x,z)$, where $x$ increases to inf\/inity and $z$ is f\/ixed
 $(z \in \BC_+)$,
is a family of embedded matrix balls such that the right semi-radii
 are uniformly bounded and the left semi-radii tend to zero. $($Recall that the $m_2\times m_1$ matrix ball, or Weyl matrix ball, with the center~$\clm$, the left semi-radius~$\clr_l$
 and the right semi-radius $\clr_r$ is the set of
 $m_2\times m_1$ matrices~$\om$
 which may may be presented
 in the form $\om = \clm + \clr_l \clu \clr_r$, where~$\clu$ are contractive $m_2\times m_1$ matrices.$)$
 \end{Remark}

 \begin{Proposition}[\cite{FKRSp1}] \label{PnW1} Let Dirac system \eqref{1.5} on $[0, \infty)$
 be given and assume that $V$ is locally summable.
 Then, the sets $\cln(x,z)$
 are well-defined. There is a unique matrix function
 $\vp(z)$ defined in $\BC_+$ and such that
\begin{gather}  \label{2.3}
\{\vp(z) \}=\bigcap_{x<\infty}\cln(x,z).
\end{gather}
This function is analytic and non-expansive $($i.e., contractive$)$.
Furthermore, this function coincides with the Weyl function of system~\eqref{1.5}.
 \end{Proposition}

Formula \eqref{2.3} is supplemented by the asymptotic relation
\begin{gather}  \label{2.3!}
\vp(z)=\lim_{b \to \infty}\vp_b(z),
\end{gather}
which is valid for any set of functions $\vp_b(z)\in \cln(b,z)$. Relation~\eqref{2.3!} follows from~\eqref{2.3}
and Remark~\ref{MB} (see also \cite[Remark~2.24]{SaSaR}).

Next, we consider the famous compatibility condition (zero curvature equation)~\eqref{zc}.
\begin{Proposition}[\cite{SaAF}] \label{TmM}
Let some $m \times m$ matrix functions $G$ and $F$ and their derivatives~$G_t$ and~$F_x$ exist on the semi-strip~$\cD$,
let $G$, $G_t$ and~$F$ be continuous with respect to~$x$ and~$t$ on~$\mathcal{D}$,
 and let~\eqref{zc}
hold. Then, we have the equality
\begin{gather} \label{2.4}
u(x,t,z)R(t,z)=R(x,t,z)u(x,0,z), \qquad R(t,z):=R(0,t,z),
\end{gather}
where $u(x,t,z)$ and $R(x,t,z)$ are normalized fundamental solutions given, respectively, by
\begin{gather} \label{2.5}
u_x=Gu, \qquad u(0,t,z)=I_m; \qquad R_t=FR, \qquad R(x,0,z)=I_m.
\end{gather}
The equality \eqref{2.4} means that the matrix function
\begin{gather*}
y(x,t,z)=u(x,t,z)R(t,z)=R(x,t,z)u(x,0,z)
\end{gather*}
satisf\/ies on $\cD$ both systems \eqref{au}. Moreover, the fundamental solution~$u$ admits the factorization
\begin{gather} \label{2.6}
u(x,t,z)=R(x,t,z)u(x,0,z)R(t,z)^{-1}.
\end{gather}
\end{Proposition}
